\documentclass[11pt]{article}

\usepackage{fontspec}
\setmainfont{times.ttf}[
  Path = ./ ,
  BoldFont = timesbd.ttf ,
  ItalicFont = timesi.ttf ,
  BoldItalicFont = timesbi.ttf ]

\usepackage{amsmath,amssymb,amsthm}
\usepackage{bm}
\usepackage[a4paper,margin=2.2cm]{geometry}
\usepackage{booktabs}
\usepackage{array}
\usepackage{enumitem}
\usepackage[dvipsnames]{xcolor}
\usepackage{titlesec}
\usepackage[hidelinks]{hyperref}
\usepackage{fancyhdr}

\definecolor{navy}{RGB}{20,40,90}
\titleformat{\section}{\normalfont\Large\bfseries\color{navy}}{\thesection.}{0.5em}{}
\titleformat{\subsection}{\normalfont\large\bfseries}{\thesubsection}{0.5em}{}

\pagestyle{fancy}\fancyhf{}
\renewcommand{\headrulewidth}{0.2pt}
\lhead{\small\itshape Phương pháp D1 — Điều khiển lai UAV LQR + NMPC-scenario chưng cất}
\cfoot{\small\thepage}

\newtheorem{lemma}{Bổ đề}
\newtheorem{proposition}{Mệnh đề}
\newtheorem{theorem}{Định lý}
\newtheorem{assumption}{Giả thiết}
\newtheorem{definition}{Định nghĩa}

\newcommand{\R}{\mathbb{R}}
\newcommand{\T}{^{\!\top}}
\newcommand{\norm}[1]{\left\lVert #1 \right\rVert}
\newcommand{\cS}{c_{S}}
\newcommand{\cL}{c_{\mathrm{LQR}}}
\newcommand{\uL}{u_{\mathrm{LQR}}}
\newcommand{\uS}{u_{\mathrm{sur}}}
\newcommand{\xref}{x_{\mathrm{ref}}}
\newcommand{\Acl}{A_{\mathrm{cl}}}

\setlength{\parskip}{0.35em}
\setlength{\parindent}{0pt}

\begin{document}

% ===================== TITLE =====================
\begin{titlepage}
\centering
\vspace*{3cm}
{\color{navy}\Huge\bfseries Phương pháp D1\par}
\vspace{0.8cm}
{\LARGE\bfseries Điều khiển lai cho UAV: LQR nền + NMPC-scenario chưng cất,\\[2mm]
với tin cậy $\cS$ (surrogate tự đoán chất lượng bám) và $\cL$ (co rút LQR)\par}
\vspace{1.2cm}
{\large\itshape Tài liệu trình bày phương pháp — bản đầy đủ để kiểm tra\par}
\vspace{2.5cm}
\begin{minipage}{0.85\textwidth}\centering
Bản ghi rõ: mô hình, tham số danh nghĩa, điều kiện bất định và nhiễu, tập quỹ đạo (họ, tốc độ,
gia tốc), thiết kế LQR (Bryson), teacher NMPC-scenario $M{=}5$, SAC tinh chỉnh $Q,R$ khởi tạo Bryson,
surrogate chưng cất, định nghĩa $\cS/\cL$, bộ điều khiển đề xuất (trộn theo tin cậy), và chứng minh
Lyapunov.
\end{minipage}
\vfill
{\itshape\small Mô phỏng; ký hiệu và tham số theo mã nguồn dự án. Nguồn tài liệu ở mục cuối.\par}
\end{titlepage}

% ===================== 0 =====================
\section{Tóm tắt phương pháp}
Mục tiêu: xây một \textbf{bộ điều khiển triển khai} cho UAV quadrotor \textbf{không cần giải tối ưu
trực tuyến}, nhưng vẫn thừa hưởng chất lượng của một bộ \textbf{NMPC-scenario} mạnh và lấy \textbf{LQR làm
nền} (phần dư chỉ bổ sung có kiểm soát; phân tích co rút là \emph{điều kiện đủ có điều kiện}, không phải
bảo chứng ổn định cho luật lai --- xem Mục 12). Kiến trúc gồm ba thành phần và một luật hợp nhất:
\begin{itemize}[leftmargin=1.4em,itemsep=1pt]
\item \textbf{LQR nền} — thiết kế theo phương pháp chuẩn (tuyến tính hoá quanh hover + Riccati rời rạc),
trọng số theo \textbf{quy tắc Bryson}. Cho ma trận Lyapunov $P$ và bộ điều khiển ổn định vùng.
\item \textbf{Teacher = NMPC-scenario ($M{=}5$)} — bộ NMPC bền vững chạy ngoại tuyến để sinh nhãn điều
khiển. \textbf{SAC} tinh chỉnh trọng số $Q,R$ của NMPC, \textbf{khởi tạo từ Bryson} để kết quả tốt ngay.
\item \textbf{Surrogate chưng cất} — mạng nơ-ron nhìn trạng thái $x$ và \emph{đồng thời đoán phần dư điều
khiển $\Delta u=u_{\mathrm{teacher}}-\uL$ và tự đoán tin cậy $\cS$} (hai đầu ra chung một mạng; chỉ học
phần teacher làm KHÁC LQR).
\item \textbf{Tin cậy $\cS,\cL$} — $\cS$ (do surrogate đoán, KHÔNG dùng $V$) là \emph{điểm chất lượng bám
vị trí gần đây} (RMS sai số 20 bước vừa qua, chuẩn hoá về $[0,1]$); $\cL$ là \emph{xác suất co rút} của
LQR (dựa trên hàm Lyapunov $V$).
\end{itemize}
Luật triển khai (không giải NMPC trực tuyến), dạng \emph{LQR $+$ phần dư có cổng}:
\[
u \;=\; \mathrm{sat}_{\mathcal U}\!\big(\uL + \alpha\,\Delta u\big), \qquad
\alpha \;=\; \cS\cdot g_L(\cL),
\]
với $\Delta u$ (phần dư teacher $-$ LQR) do surrogate đoán, $\mathrm{sat}_{\mathcal U}$ là chặn biên
actuator, và $g_L$ là \textbf{cổng do tin cậy LQR điều khiển} (Mục 11). \textbf{LQR được ưu tiên}: khi LQR
tự tin thì $g_L\to0$ nên $u\approx\uL$; residual chỉ hoạt động mạnh khi \emph{LQR kém tin ($g_L$ lớn) VÀ
surrogate đáng tin ($\cS$ cao)}. Nếu cả hai cùng cao $\Rightarrow$ vẫn thuần LQR.

% ===================== 1 =====================
\section{Kiến trúc tổng thể}
Một lần huấn luyện cho mỗi hạt giống (seed) gồm \textbf{hai khối offline}. Khối A (SAC $+$ chưng cất
residual) chạy \emph{đồng thời} --- KHÔNG đóng băng $Q,R$: surrogate chỉ nhìn trạng thái $x$ (không nhìn
$Q,R$) nên học được ánh xạ $x\to\Delta u$ qua mọi $Q,R$ mà SAC thử, \emph{kể cả trường hợp xấu} (nhờ
replay ưu tiên nhãn gần đây, cuối cùng hội tụ về teacher ứng với $Q,R$ mà SAC chốt); những vùng surrogate
bám kém sẽ được $\cS$ tự hạ thấp. Khối B (huấn luyện $\cS$) tách ra CHỈ để tránh vòng phụ thuộc
$\alpha\leftrightarrow\cS$.
\begin{center}\small
\begin{tabular}{@{}p{3.4cm} p{4.2cm} p{6.6cm}@{}}
\toprule
\textbf{Khối} & \textbf{Bộ điều khiển bay} & \textbf{Dữ liệu thu}\\
\midrule
A. SAC $+$ residual (đồng thời) & NMPC teacher (SAC tune $Q,R$) $+$ LQR song song &
$Q,R$; nhãn phần dư $\Delta u^*=u_T-\uL$ (đầu $\Delta u$); nhãn co rút LQR (cho $\cL$)\\
B. Huấn luyện $\cS$ & LQR $+$ residual NN, \textbf{$\alpha=1$} (rollout surrogate) &
RMS sai số bám 20 bước vừa qua $\Rightarrow$ nhãn $s_k$ (đầu $\cS$)\\
\bottomrule
\end{tabular}
\end{center}
\textbf{Tránh vòng phụ thuộc.} Ở Khối B, $\cS$ chưa đáng tin khi mới học, nên \emph{không} dùng ngay $\cS$
để quyết $\alpha$ rồi lấy kết quả đó dạy lại $\cS$. Thay vào đó cho nhánh residual chạy với \textbf{luật
cố định $\alpha=1$}: $u_k=\mathrm{sat}_{\mathcal U}(\uL+\widehat{\Delta u}_k)$, thu $(z_k,s_k)$ để huấn
luyện đầu $\cS$ (đóng băng thân mạng $+$ đầu $\Delta u$).

Toàn bộ trạng thái lưu \textbf{checkpoint}. Khi triển khai chỉ còn LQR (rẻ) $+$ surrogate (một lần suy
luận) $+$ hai tin cậy; \textbf{không} giải NMPC --- chi phí ở mức LQR, chất lượng hướng tới NMPC.

% ===================== 2 =====================
\section{Mô hình quadrotor 12 trạng thái}
Trạng thái $x=[\,p\,(3);\ \eta\,(3);\ v\,(3);\ \omega\,(3)\,]\in\R^{12}$ gồm vị trí $p$, góc Euler
$\eta=(\phi,\theta,\psi)$, vận tốc $v$, tốc độ góc thân $\omega$. Đầu vào
$u=[\,T;\ \tau_\phi;\ \tau_\theta;\ \tau_\psi\,]$ (lực đẩy tổng và ba mô-men). Khung quán tính $z$ hướng
lên; hover tại góc $0$ dùng $T=mg$. Động lực học phi tuyến 6-DOF (Newton–Euler) tích phân bằng RK4, bước
$T_s$.

\begin{center}\small
\begin{tabular}{@{}p{5.2cm} c p{6.2cm} c@{}}
\toprule
\textbf{Tham số} & \textbf{Ký hiệu} & \textbf{Giá trị} & \textbf{Đơn vị}\\
\midrule
Khối lượng & $m$ & $0.486$ & kg\\
Gia tốc trọng trường & $g$ & $9.81$ & m/s$^2$\\
Quán tính (đường chéo) & $J$ & $\mathrm{diag}(3.828,\,3.828,\,7.657)\times 10^{-3}$ & kg$\cdot$m$^2$\\
Cánh tay đòn & $\ell$ & $0.25$ & m\\
Cản tịnh tiến & $D_v$ & $[5.567,\,5.567,\,6.354]\times 10^{-4}$ & N$\cdot$s/m\\
Cản quay & $D_\omega$ & $[5.567,\,5.567,\,6.354]\times 10^{-4}$ & N$\cdot$m$\cdot$s\\
Hệ số lực/mô-men & $\alpha_T,\alpha_\tau$ & $1.0\ ;\ [1,1,1]$ & --\\
Giới hạn lực đẩy & $T$ & $[0,\ 4mg]=[0,\ 19.07]$ & N\\
Giới hạn mô-men & $\tau$ & $\pm0.5\,;\ \pm0.5\,;\ \pm0.25$ & N$\cdot$m\\
\bottomrule
\end{tabular}
\end{center}
{\small\itshape Nguồn: step1\_plant\_config.}

% ===================== 3 =====================
\section{Điều kiện thử nghiệm: bất định mô hình và nhiễu}
Để bài toán có ý nghĩa bền vững, mỗi lần bay dùng \textbf{plant thật lệch danh nghĩa} (sai số mô hình) và
\textbf{nhiễu ngoại lực/mô-men}. Teacher NMPC hedge sai số này bằng $M{=}5$ kịch bản; LQR chỉ thiết kế
trên mô hình danh nghĩa (không biết trước lệch).

\begin{center}\small
\begin{tabular}{@{}p{5.6cm} c c@{}}
\toprule
\textbf{Nhóm tham số} & \textbf{train $\rho$ (in-dist.)} & \textbf{ood $\rho$}\\
\midrule
Khối lượng $m$ & $\pm10\%$ & $\pm25\%$\\
Quán tính $I_x,I_y,I_z$ & $\pm5\%$ & $\pm15\%$\\
Cản tịnh tiến $D_v$ & $\pm30\%$ & $\pm60\%$\\
Cản quay $D_\omega$ & $\pm30\%$ & $\pm60\%$\\
Hệ số lực đẩy $\alpha_T$ & $\pm15\%$ & $\pm30\%$\\
Hệ số mô-men $\alpha_\tau$ & $\pm15\%$ & $\pm30\%$\\
\bottomrule
\end{tabular}\qquad
\begin{tabular}{@{}p{2.6cm} c c@{}}
\toprule
\textbf{Nhiễu sine} & \textbf{Biên độ} & \textbf{Tần (Hz)}\\
\midrule
Lực $(x,y,z)$ & $[0.05,0.05,0.08]$ N & $[0.7,0.5,0.3]$\\
Mô-men $(\phi,\theta,\psi)$ & $[2,2,1]\times10^{-3}$ N$\cdot$m & $[0.9,0.6,0.4]$\\
\bottomrule
\end{tabular}
\end{center}
{\small\itshape Nguồn: step1\_plant\_config.uncertainty và .disturbance.training.}

% ===================== 4 =====================
\section{Tập quỹ đạo tham chiếu và cấu trúc episode}
Tập quỹ đạo có \textbf{5 họ hình học}, quét theo \textbf{tốc độ đỉnh} và \textbf{gia tốc mục tiêu}. Mỗi
quỹ đạo dài \textbf{1000 bước $=50$ s} ($T_s=0.05$ s).

\begin{center}\small
\begin{tabular}{@{}p{4.6cm} p{9.8cm}@{}}
\toprule
\textbf{Thuộc tính} & \textbf{Giá trị (nguồn: targeted\_lqr\_weak\_config)}\\
\midrule
5 họ quỹ đạo & circle, lemniscate, vertical\_circle, spatial\_helix, smooth\_waypoints\\
Tốc độ đỉnh (8 mốc, in-dist.) & $0.5,\ 1,\ 2,\ 4,\ 6,\ 8,\ 10,\ 12$ (m/s)\\
Tốc độ ngoài phân phối & $14,\ 16$ (m/s)\\
Gia tốc mục tiêu (3 mức) & $2.0$ (mild), $5.0$ (moderate), $9.0$ (hard) (m/s$^2$)\\
Độ cao / jitter hình học & $[1,3]$ m $/\ \times[0.90,1.10]$\\
Độ dài / bước & 1000 bước $=50$ s, $T_s=0.05$ s\\
\bottomrule
\end{tabular}
\end{center}

\subsection{Episode và điều kiện công bằng LQR vs teacher}
\textbf{1 episode (1 lần chấm SAC) $=$ chạy cả 5 họ quỹ đạo, mỗi họ 1000 bước}, trên plant thật có nhiễu
$+$ sai số mô hình. \textbf{LQR chạy đúng cùng tập quỹ đạo và cùng điều kiện} với teacher, nhưng
\textbf{độc lập từng quỹ đạo} (không hedge kịch bản như teacher) và \textbf{thiết kế vẫn trên mô hình
danh nghĩa} — vì 5 bối cảnh lệch là không biết trước. Chia train/test theo băm (hash); tập test giữ
nguyên.

% ===================== 5 =====================
\section{Bộ điều khiển nền LQR}
\subsection{Phương pháp}
Phương pháp: \textbf{LQR vô hạn thời gian, rời rạc, trên mô hình tuyến tính hoá Jacobian quanh điểm cân
bằng hover}, giải bằng \textbf{phương trình Riccati đại số rời rạc (DARE)}.
\begin{enumerate}[leftmargin=1.6em,itemsep=1pt]
\item Điểm cân bằng hover: $x_h=0_{12}$, $u_h=[mg;0;0;0]$ (kiểm chứng $\dot x=0$).
\item Tuyến tính hoá Jacobian số quanh $(x_h,u_h)$ $\rightarrow$ $(A,B)$.
\item Rời rạc hoá giữ bậc 0 (ZOH) với $T_s=0.05$ s $\rightarrow$ $(A_d,B_d)$.
\item Giải DARE với $(A_d,B_d,Q,R)$ $\rightarrow$ độ lợi $K$ và nghiệm Riccati $S$.
\end{enumerate}
\[
u \;=\; u_h - K\,(x-\xref).
\]
Ma trận $\boxed{P=S}$ (nghiệm Riccati) chính là ma trận Lyapunov dùng cho tiêu chí co rút (Mục 9, 11);
không tính lại, không dùng $Q_f$.

\subsection{Chọn trọng số — quy tắc Bryson}
Trọng số theo \textbf{quy tắc Bryson}: mỗi phần tử là nghịch đảo bình phương độ lệch tối đa chấp nhận
được:
\[
Q_{ii}=\frac{1}{(e_{\mathrm{allow},i})^2},\qquad
R_{jj}=\frac{1}{(du_{\mathrm{allow},j})^2}.
\]
\begin{center}\small
\begin{tabular}{@{}p{4.6cm} c c@{}}
\toprule
\textbf{Kênh trạng thái} & \textbf{$e_{\mathrm{allow}}$} & \textbf{$Q_{ii}$}\\
\midrule
Vị trí $p_x,p_y,p_z$ (m) & $0.10$ & $100.0$\\
Góc $\phi,\theta,\psi$ (rad) & $5^\circ=0.0873$ & $131.3$\\
Vận tốc $v_x,v_y,v_z$ (m/s) & $0.30$ & $11.11$\\
Tốc độ góc $p,q,r$ (rad/s) & $2.0$ & $0.25$\\
\bottomrule
\end{tabular}\qquad
\begin{tabular}{@{}p{3.0cm} c c@{}}
\toprule
\textbf{Kênh điều khiển} & \textbf{$du_{\mathrm{allow}}$} & \textbf{$R_{jj}$}\\
\midrule
Lực đẩy $T$ (N) & $mg=4.768$ & $0.0440$\\
Mô-men $\tau_\phi$ (N$\cdot$m) & $0.5$ & $4.0$\\
Mô-men $\tau_\theta$ (N$\cdot$m) & $0.5$ & $4.0$\\
Mô-men $\tau_\psi$ (N$\cdot$m) & $0.25$ & $16.0$\\
\bottomrule
\end{tabular}
\end{center}

\noindent\textbf{Minh bạch:} quy tắc Bryson được trích nguồn, nhưng các giá trị $e_{\mathrm{allow}}$ (vị
trí $0.1$ m, tư thế $5^\circ$, vận tốc $0.3$ m/s, tốc độ góc $2$ rad/s) là \emph{dung sai thiết kế} do
người thiết kế đặt, đúng như Bryson \& Ho yêu cầu chọn ``độ lệch tối đa chấp nhận được''. LQR dùng
feedforward hover $u_h$ (không có feedforward theo quỹ đạo) nên là bộ điều chỉnh quanh hover — hạn chế
đã biết khi bay xa hover.

\noindent{\small\itshape Nguồn: Bouabdallah, Noth, Siegwart (IROS 2004) — LQ quadrotor tuyến tính hoá
quanh hover; Bryson \& Ho (1975) — quy tắc Bryson; Sir Elkhatem \& Engin (2022) — Bryson cho LQR trên
UAV; Anderson \& Moore (1990) — LQR/DARE.}

% ===================== 6 =====================
\section{Teacher: NMPC-scenario bền vững ($M=5$)}
Teacher tối ưu \textbf{một chuỗi điều khiển duy nhất} sao cho tốt cho cả $M{=}5$ kịch bản plant bất định.
Trạng thái tăng cường (64): $X=[\,x^{(1)}\!,\dots,x^{(5)}\,(60);\ u_{\mathrm{prev}}\,(4)\,]$; đầu vào
$\delta u$; điều khiển áp dụng $u=u_{\mathrm{prev}}+\delta u$ (delta-u). Mỗi kịch bản $i$ tích phân RK4
với bộ tham số lệch $\theta_i$.
\[
\min\ \ \frac{1}{M}\sum_{i=1}^{M}\sum_{k=0}^{N-1}
\Big(\norm{x^{(i)}_k-x_{\mathrm{ref},k}}_{Q}^2 + \norm{u_k-u_h}_{R}^2\Big).
\]
\begin{itemize}[leftmargin=1.4em,itemsep=1pt]
\item Chân trời dự báo $N=20$; chân trời điều khiển $N_c=5$ (move-blocking: $\delta u=0$ cho các stage
$N_c,\dots,N-1$).
\item Ràng buộc actuator áp qua chặn trạng thái $u_{\mathrm{prev}}\in[u_{\min},u_{\max}]$.
\item Trọng số $Q,R$ \textbf{khởi tạo Bryson}, sau đó do SAC tinh chỉnh (Mục 7); cập nhật nóng
\texttt{cost\_W} không dựng lại solver.
\item Giải bằng \textbf{acados} (SQP thực $+$ tìm kiếm đường merit, tích phân DISCRETE, QP HPIPM), mô
hình ký hiệu \textbf{CasADi}. $N=20,\ N_c=5,\ M=5$ là tham số cố định.
\end{itemize}
{\small\itshape Nguồn: scenario-MPC — Calafiore \& Campi (2006); acados — Verschueren et al. (2022);
CasADi — Andersson et al. (2019).}

\noindent{\small\itshape Giới hạn đã biết: trung bình chi phí trên $M=5$ kịch bản KHÔNG tự tạo bảo đảm
robust trên toàn miền bất định; đây là hedge kinh nghiệm, không phải chứng nhận min--max.}

% ===================== 7 =====================
\section{SAC tinh chỉnh $Q,R$ (khởi tạo Bryson)}
Đây là \textbf{Khối A} (Mục kiến trúc). \textbf{Ý đồ:} khởi tạo $Q,R=$ Bryson để reward \textbf{tốt
ngay từ đầu}; SAC chỉ tinh chỉnh quanh đó. Song song, quỹ đạo \textbf{LQR} thu ở đây cho \textbf{nhãn co
rút của $\cL$}. (Nhãn residual của teacher lấy cùng ở Khối A; nhãn $\cS$ lấy ở Khối B từ rollout
surrogate --- KHÔNG lấy ở đây.)
\begin{itemize}[leftmargin=1.4em,itemsep=1pt]
\item Hành động $a\in[-1,1]^6$ ánh xạ sang \textbf{log-multiplier} cho 4 nhóm $Q$ (pos/att/vel/rate) và 2
nhóm $R$ ($T/\tau$), dải $[10^{-1.5},10^{1.5}]$ \textbf{quanh Bryson}; $a=0\Rightarrow$ đúng Bryson.
\item SAC dạng \textbf{bandit 1 bước} ($\gamma=0$): actor Gaussian trạng thái-độc-lập $(\mu,\log\sigma)$,
hai critic $Q(a)$, hệ số entropy $\alpha_{\mathrm{ent}}$ tự chỉnh; cập nhật Adam.
\item Phần thưởng episode:
$\ r=-\big(\mathrm{RMSE}_{\mathrm{pos}} + 0.01\sqrt{\textstyle\sum\norm{\Delta u}^2}
+ 0.5\,r_{\mathrm{viol}} + 5\,r_{\mathrm{fail}}\big)$ (RMSE chặn $5$ m/bước).
\item \textbf{Checkpoint đầy đủ} (actor/critic/target/$\alpha_{\mathrm{ent}}$/optimizer/replay $+$
surrogate $+$ RNG $+$ $Q,R$) để \textbf{resume}.
\end{itemize}
{\small\itshape Nguồn: Haarnoja et al. (2018), Soft Actor-Critic. Lưu ý: dùng RL để tinh chỉnh trọng số
MPC cố định có ít tài liệu nền; Bryson chỉ đóng vai KHỞI TẠO cho SAC. Kết quả SAC tốt trên tập train
KHÔNG bảo đảm tổng quát hoá --- cần báo reward, ngân sách đánh giá và hiệu năng trên \emph{tập test tách
biệt}.}

% ===================== 8 =====================
\section{Surrogate chưng cất (behavioral cloning theo dòng)}
\begin{itemize}[leftmargin=1.4em,itemsep=1pt]
\item \textbf{Đặc trưng 208 chiều}: lịch sử trạng thái ($12{\times}4$) $+$ lịch sử đầu vào ($4{\times}4$)
$+$ nhìn trước tham chiếu ($12{\times}11$) $+$ phần dư ($12$).
\item \textbf{Mạng hai đầu ra}: ký hiệu $z_k$ là đầu vào NN 208 chiều tại bước $k$. Thân chung MLP
$208$-$128$-$128$-$128$ (swish) rẽ hai nhánh: \emph{nhánh phần dư} $\to 4$ cho $\widehat{\Delta u}$, và
\emph{nhánh tin cậy} $\to 1$ (sigmoid) cho $\cS$.
\item \textbf{Chuẩn hoá} cố định theo thang vật lý; huấn luyện \textbf{streaming} (adamupdate, minibatch
$256$), \textbf{replay ưu tiên nhãn gần đây}.
\item \textbf{Nhánh $\Delta u$ --- nhãn từ TEACHER (Khối A)}: nhãn phần dư
$\Delta u^*_k=u_{\mathrm{teacher},k}-\uL(x_k)$ với $\uL(x_k)=u_h-K(x_k-\xref)$ (LQR tính tại \emph{cùng}
trạng thái); chỉ nạp mẫu khi teacher cho nghiệm hữu hạn. Đây là ``chưng cất NMPC'' dạng phần dư.
\item \textbf{Nhánh $\cS$ --- nhãn từ ROLLOUT SURROGATE (Khối B)}: học điểm $s_k=$ RMS sai số bám vị
trí 20 bước vừa qua của rollout surrogate-residual offline (định nghĩa ở Mục 10.2).
\end{itemize}
Hai đầu \textbf{dùng chung thân mạng} nhưng \textbf{khác nguồn nhãn}. Mất mát chung:
\[
\mathcal{L}=\lambda_u\,\norm{\widehat{\Delta u}_k-\Delta u^*_k}^2+\lambda_c\,\big(\cS(z_k)-s_k\big)^2,
\qquad \widehat{\Delta u}=\pi_\theta^{u}(z_k),\ \ \cS=\pi_\theta^{c}(z_k).
\]
\textbf{Lưu ý dữ liệu}: nhãn $\Delta u^*_k$ chỉ có khi teacher được tính tại đúng $x_k$. Nên ở Khối B,
muốn cập nhật cả đầu residual thì phải \emph{gọi teacher offline tại các trạng thái surrogate đã đi qua};
nếu không, chỉ cập nhật đầu $\cS$ và giữ đầu residual cố định.
{\small\itshape Behavioral cloning chịu ``dịch phân phối''; hướng cải thiện mạnh nhất là DAgger (Ross et
al., 2011). Đầu $\cS$ là cơ chế để surrogate \emph{tự biết} vùng nó bám kém.}

% ===================== 9 =====================
\section{Định nghĩa tin cậy $\cS$ và $\cL$}
Hai tin cậy được định nghĩa theo \textbf{tiêu chí phù hợp với từng bộ điều khiển}, KHÔNG dùng chung một
hàm $V$: $V$ là hàm Lyapunov \emph{của LQR}, nên dùng nó để chấm surrogate là không hợp lí.

\subsection{$\cL$ --- dựa trên hàm Lyapunov của LQR (giữ nguyên)}
$V(e)=e\T P e$ với $P=S$ (Riccati của LQR), $e=x-\xref$, là hàm Lyapunov tự nhiên của LQR $\Rightarrow$ đo
co rút của LQR bằng $V$ là hợp lí. Trên cửa sổ $H=20$ bước:
\[
\Delta V_H(k)=V(k{+}H)-V(k),\qquad y_L(k)=\mathbf{1}\!\left[\Delta V_H(k)<0\right],
\]
\[
\cL(e)=\Pr\!\big(y_L=1\mid e\big)=\sigma\!\big(w\T z(e)+b\big)\quad(\text{hồi quy logistic trên LQR vòng
kín}),
\]
với $z(e)$ chuẩn hoá của $\varphi(e)=[\,e\,(12);\ \norm{e_{\mathrm{pos}}},\norm{e_{\mathrm{att}}},
\norm{e_{\mathrm{vel}}},\norm{e_{\mathrm{rate}}}\,]\in\R^{16}$.

\subsection{$\cS$ --- điểm chất lượng bám vị trí gần đây, do surrogate đoán (KHÔNG dùng $V$, KHÔNG so LQR)}
Trên rollout surrogate, sai số bám vị trí $e^S_{p,k}=p^S_k-p_{\mathrm{ref},k}$. Lấy \textbf{RMS vị trí trên
$20$ bước VỪA QUA}:
\[
E^S_{20}(k)=\sqrt{\frac{1}{20}\sum_{j=k-19}^{k}\norm{e^S_{p,j}}^2}\,.
\]
Chuẩn hoá RMS về $[0,1]$ với thang $\varepsilon_p$ (do người thiết kế chọn và công bố), ví dụ:
\[
s_k=\exp\!\left(-\Big(\frac{E^S_{20}(k)}{\varepsilon_p}\Big)^2\right)\in[0,1].
\]
RMS gần $0\Rightarrow s_k\to1$; RMS lớn $\Rightarrow s_k\to0$. \emph{UAV bám hoàn hảo ($e\equiv0$) suốt $20$
bước nhận $s_k=1$}, không còn bị phạt vô lí như kiểu ``đếm số bước giảm''. (Hàm mũ chỉ là một lựa chọn
chuẩn hoá; có thể thay bằng phép khác, miễn $\cS$ học \emph{mức} sai số bám, không phải tần suất giảm.)

\begin{definition}[Tin cậy $\cS$ --- điểm chất lượng bám gần đây]
Đầu $\cS$ của surrogate học dự đoán $s_k$ từ đầu vào NN hiện tại $z_k$:
\[
\cS(z_k)=\pi_\theta^{c}(z_k)\ \approx\ s_k.
\]
Vì cửa sổ là \textbf{20 bước vừa qua}, $\cS$ là \emph{điểm dự đoán chất lượng bám vị trí gần đây}, KHÔNG
phải ``xác suất sai số giảm trong 20 bước tới''. Khi bay chỉ cần một lần suy luận từ $z_k$; \textbf{không
tính cửa sổ 20 bước online}.
\end{definition}

Nhãn $s_k$ lấy từ \textbf{rollout surrogate-residual offline} (Khối B). Diễn giải: $\cS$ cao
$\Rightarrow$ gần đây surrogate bám tốt $\Rightarrow$ đáng tin, cho residual chạy; $\cS$ thấp $\Rightarrow$
bám kém $\Rightarrow$ kém tin. Cả $\cS,\cL\in[0,1]$.

{\small\itshape Nguồn: $V$ = nghiệm Riccati là hàm Lyapunov của LQR (Anderson \& Moore, 1990); hồi quy
logistic (Bishop, 2006).}

\subsection{Ý nghĩa thống kê --- những điều $\cS,\cL$ KHÔNG bảo đảm}
\begin{itemize}[leftmargin=1.4em,itemsep=1pt]
\item \textbf{$\cS$ đo chất lượng bám gần đây}, không đồng nghĩa với: (i) surrogate mô phỏng đúng teacher
(sai số chưng cất nhỏ), hay (ii) residual bảo toàn co rút Lyapunov. Ba thuộc tính này \emph{không} tương
đương --- surrogate có thể bám tốt mà vẫn khác teacher, hoặc mô phỏng teacher đúng ở nơi cả hai đều sai số
lớn do giới hạn vật lý. Vì vậy KHÔNG diễn giải $\cS$ là ``xác suất surrogate chính xác'' hay ``xác suất
residual an toàn''.
\item \textbf{Dịch phân phối:} nhãn $\cS$ thu ở $\alpha=1$ (residual đầy đủ), còn triển khai
$\alpha=\cS g_L(\cL)$ tạo phân phối trạng thái khác. Cần \emph{kiểm định $\cS$ trên rollout của chính bộ
điều khiển lai}, không chỉ rollout residual đầy đủ.
\item \textbf{$\cL$ là xác suất thống kê}, không phải chứng nhận co rút: $\cL(e)=\Pr(V_{k+20}-V_k<0)$ cao
KHÔNG bảo đảm $V_{k+1}-V_k<0$ trên plant hiện tại (nhất là khi quỹ đạo động, plant lệch, hay bão hoà). Hơn
nữa nhãn $\cL$ dùng co rút \emph{sau 20 bước} còn phân tích Mục 12 dùng co rút \emph{từng bước} --- hai đại
lượng khác nhau, không thay thế nhau.
\end{itemize}

% ===================== 10 =====================
\section{Bộ điều khiển đề xuất (blend triển khai)}
Không giải NMPC trực tuyến. LQR là ``xương sống'' và được \textbf{ưu tiên}; surrogate chỉ bổ sung
\textbf{phần dư} $\Delta u=\pi_S^u(x)$ \emph{khi được phép}. Định nghĩa \textbf{cổng do LQR điều khiển}
(theo tin cậy $\cL$):
\[
g_L(\cL)=\mathrm{clip}\!\left(\frac{c_{\mathrm{high}}-\cL}{\,c_{\mathrm{high}}-c_{\mathrm{low}}\,},\,0,\,1\right),
\qquad 0\le c_{\mathrm{low}}<c_{\mathrm{high}}\le 1.
\]
Trọng số residual và luật điều khiển:
\[
\boxed{\,u \;=\; \mathrm{sat}_{\mathcal U}\!\big(\uL + \alpha\,\Delta u\big),\qquad
\alpha \;=\; \underbrace{\cS}_{\substack{\text{surrogate}\\\text{đáng dùng?}}}\;\cdot\;
\underbrace{g_L(\cL)}_{\substack{\text{LQR có cho}\\\text{residual chạy?}}}\,},\qquad \alpha\in[0,1].
\]
Hành vi (đúng như thiết kế):
\begin{itemize}[leftmargin=1.4em,itemsep=1pt]
\item $\cL\ge c_{\mathrm{high}}$ (LQR rất tự tin) $\Rightarrow g_L=0\Rightarrow \alpha=0\Rightarrow u=\uL$:
\textbf{thuần LQR}, bất kể $\cS$ (cả hai cùng cao vẫn thuần LQR --- LQR ưu tiên).
\item $\cL\le c_{\mathrm{low}}$ (LQR kém tin) và $\cS$ cao $\Rightarrow g_L=1,\ \alpha=\cS$ lớn
$\Rightarrow$ \textbf{residual hoạt động mạnh} (surrogate sửa cho LQR); $\cS=1\Rightarrow u=u_{\mathrm{teacher}}$.
\item $\cL$ thấp nhưng $\cS$ thấp $\Rightarrow \alpha$ nhỏ $\Rightarrow$ vẫn gần LQR (không tin surrogate).
\end{itemize}
Ý nghĩa: dùng LQR bất cứ khi nào nó đủ tin; chỉ nhường cho residual ở nơi \emph{LQR không chắc mà
surrogate lại chắc}. $c_{\mathrm{low}},c_{\mathrm{high}}$ là ngưỡng thiết kế.

\noindent\emph{Lưu ý:} đây là \textbf{cơ chế lựa chọn điều khiển dựa trên dữ liệu}, KHÔNG phải cơ chế
chứng nhận Lyapunov. Luật này không tự bảo đảm $\alpha\le\bar\alpha(e,d)$ (ngân sách co rút ở Mục 12), nên
tính ổn định của bộ điều khiển lai chỉ có ở dạng \emph{điều kiện đủ có điều kiện} (Mục 12), trừ khi bổ
sung cơ chế giới hạn $\alpha_{\mathrm{safe}}$ (Mệnh đề~\ref{prop2}).

% ===================== 11 =====================
\section{Phân tích co rút (điều kiện đủ, KHÔNG phải bảo chứng cho luật lai)}
\textbf{Phạm vi (nêu thẳng).} Phần này chỉ chứng minh một \emph{điều kiện đủ} để phần dư $\alpha\Delta u$
bảo toàn tính co rút của $V(e)=e\T P e$, dưới các giả thiết hẹp; nó \textbf{không} khẳng định bộ điều khiển
lai $u=\mathrm{sat}_{\mathcal U}(\uL+\alpha\Delta u)$, $\alpha=\cS g_L(\cL)$, luôn ổn định.

\begin{assumption}[Phạm vi hẹp]\label{as}
(A1) dùng mô hình \emph{tuyến tính hoá} quanh hover; (A2) actuator \emph{chưa bão hoà}; (A3) tham chiếu
\emph{khả thi} sao cho feedforward $w_k=A_d r_k-r_{k+1}=0$ (bài toán \emph{điều chỉnh} quanh hover, KHÔNG
phải bám quỹ đạo động tổng quát); (A4) \emph{không nhiễu, không sai số mô hình}. Ngoài ra
$P=S=S\T\succ0$ từ DARE của $(A_d,B_d,Q,R)$, $Q,R\succ0$, $\Acl=A_d-B_dK$ Schur.
\end{assumption}

\begin{lemma}[LQR là Lyapunov, dưới A1--A4]\label{lem1}
Dưới A1,A3,A4 sai số kín LQR là $e_{k+1}=\Acl e_k$ (vì $w_k=0$). Đồng nhất thức Riccati
$\Acl\T S\Acl-S=-(Q+K\T RK)$ cho $\ V(e_{k+1})-V(e_k)=-\,e_k\T(Q+K\T RK)e_k<0$ với $e_k\ne0$.
\end{lemma}
\begin{proof}
Thay $e_{k+1}=\Acl e_k$ vào $V$ và dùng đồng nhất thức Riccati; $Q+K\T RK\succ0$.
\end{proof}
\noindent Dưới A2--A4, kết quả mở rộng \emph{cục bộ} quanh hover cho plant phi tuyến bằng phương pháp gián
tiếp Lyapunov (chỉ cho bài toán điều chỉnh, $w=0$).

\subsection{Điều kiện đủ cho trọng số residual (đã sửa lỗi $q=0$)}
Dưới A1--A4, đặt phần dư $d=\Delta u$; blend chưa bão hoà $u=\uL+\alpha d$ cho $e_{k+1}=\Acl e+\alpha B_d d$,
nên
\begin{equation}
\Delta V_{\mathrm{bl}}
= \underbrace{-\,e\T(Q+K\T RK)e}_{-\,m(e)\ <\ 0}
+\alpha\underbrace{2\,e\T\Acl\T S B_d d}_{b(e,d)}
+\alpha^2\underbrace{d\T B_d\T S B_d d}_{q(d)\ \ge\ 0}. \label{eq:dv}
\end{equation}

\begin{proposition}[Ngân sách co rút $\bar\alpha$]\label{prop}
$\Delta V_{\mathrm{bl}}<0$ với mọi $\alpha\in[0,\bar\alpha(e,d))$, trong đó
\[
\bar\alpha(e,d)=
\begin{cases}
\dfrac{-b+\sqrt{b^2+4qm}}{2q}, & q>0,\\[2mm]
m/b, & q=0,\ b>0,\\[1mm]
+\infty, & q=0,\ b\le0.
\end{cases}
\]
Hơn nữa, vì $S\succ0$ nên $q=d\T B_d\T S B_d d=0\Rightarrow B_d d=0\Rightarrow b=2e\T\Acl\T S B_d d=0$; do
đó với cấu trúc bài toán này nhánh $q=0$ luôn cho $\bar\alpha=+\infty$ (nhánh $q=0,b>0$ không xảy ra).
\end{proposition}
\begin{proof}
$h(\alpha)=-m+b\alpha+q\alpha^2$, $h(0)=-m<0$. Với $q>0$, $h$ có đúng một nghiệm dương $\bar\alpha$, $h<0$
trên $[0,\bar\alpha)$. Với $q=0$: $h=-m+b\alpha$ âm khi $\alpha<m/b$ nếu $b>0$, và âm mọi $\alpha\ge0$ nếu
$b\le0$.
\end{proof}

\subsection{Khoảng trống với luật triển khai}
Luật thực tế $\alpha=\cS\,g_L(\cL)$ chỉ bảo đảm $\alpha\in[0,1]$; \textbf{không} có ràng buộc nào ép
$\alpha\le\bar\alpha(e,d)$. Vì vậy Mệnh đề~\ref{prop} \textbf{không} kéo theo $\Delta V<0$ cho luật lai:
hai cổng $\cS,g_L$ là \emph{cơ chế lựa chọn dựa trên dữ liệu}, KHÔNG phải chứng nhận Lyapunov. Tài liệu do
đó \textbf{không} khẳng định bộ điều khiển lai ổn định.

\begin{proposition}[Cơ chế bảo đảm tuỳ chọn, trong phạm vi A1--A4]\label{prop2}
Thay trọng số bằng $\ \alpha_{\mathrm{safe}}=\min\!\big\{\cS\,g_L(\cL),\,(1-\varepsilon)\,\bar\alpha(e,d)\big\}$,
$0<\varepsilon<1$, thì $\alpha_{\mathrm{safe}}<\bar\alpha\Rightarrow\Delta V_{\mathrm{bl}}<0$ dưới A1--A4 ---
đóng khoảng trống \emph{trong phạm vi hẹp}. Đây là điều kiện đủ có điều kiện, CHƯA xử lý (A2) bão hoà,
(A3) tham chiếu động, (A4) nhiễu/sai số mô hình.
\end{proposition}

\subsection{Ngoài phạm vi (chưa chứng minh --- không đổi tên kết luận)}
\begin{itemize}[leftmargin=1.4em,itemsep=1pt]
\item \textbf{Bám quỹ đạo động ($w\ne0$):} mô hình tuyến tính hoá quanh hover có đầu vào là \emph{độ lệch}
$(u_k-u_h)$, tức $x_{k+1}=A_d x_k+B_d(u_k-u_h)$; thay $u_k=u_h-Ke_k$ và $x_k=e_k+r_k$ được
$e_{k+1}=\Acl e_k+w_k$ với $\boxed{w_k=A_d r_k-r_{k+1}}$ (không có số hạng $B_d u_h$). Với quỹ đạo động,
$r_{k+1}\ne A_d r_k$ nên $w_k\ne0$: LQR chỉ có feedforward hover $u_h$, không có feedforward theo gia tốc
quỹ đạo, nên dù $e_k=0$ vẫn có thể $e_{k+1}\ne0$. Cần đưa $w$ vào và chứng minh \emph{chặn cuối} (ISS)
dạng $\Delta V\le-c\norm{e}^2+\gamma\norm{w}^2$ --- \emph{chưa} làm ở đây. (Trường hợp $x_h\ne0$ tổng
quát: $w_k=x_h+A_d(r_k-x_h)-r_{k+1}$.)
\item \textbf{Bão hoà actuator:} $u=\mathrm{sat}_{\mathcal U}(\uL+\alpha d)$; kể cả $\alpha=0$ cũng chỉ có
$u=\mathrm{sat}_{\mathcal U}(\uL)\ne\uL$ nói chung. Bảo chứng chỉ còn khi LQR chưa bão hoà.
\item \textbf{Nhiễu, sai số mô hình:} dưới nhiễu không triệt tiêu không thể suy $\Delta V<0$ mọi trạng thái
$\ne0$; mục tiêu đúng là ổn định thực dụng/chặn cuối, cần chứng minh riêng và xác định miền áp dụng.
\end{itemize}
\noindent\textbf{Kết luận (thu hẹp).} Kết quả lý thuyết ở đây là \emph{điều kiện đủ để phần dư bảo toàn co
rút trên mô hình tuyến tính hoá, chưa bão hoà, không nhiễu, tham chiếu cân bằng} --- KHÔNG phải định lý
khẳng định bộ điều khiển lai luôn ổn định. Nền lý thuyết: Anderson \& Moore (1990); Khalil, \emph{Nonlinear
Systems} (2002).

% ===================== 12 =====================
\section{Chỉ số đánh giá}
\begin{itemize}[leftmargin=1.4em,itemsep=1pt]
\item \textbf{Bám quỹ đạo}: RMSE và max sai số vị trí (căn lề bước $k$ so với $x_{\mathrm{ref},k+1}$).
\item \textbf{Co rút}: fractionContracting $=$ tỉ lệ cửa sổ $H$ có $\Delta V_H<0$; $g_H$ trung bình.
\item \textbf{Tin cậy}: độ chính xác phân loại $\cS,\cL$ và tỉ lệ nền co rút.
\item \textbf{So sánh}: LQR thuần vs teacher NMPC vs surrogate thuần vs blend, trên plant danh nghĩa và
\textbf{off-nominal} (cùng plant lệch cho mọi bộ điều khiển).
\end{itemize}

% ===================== 13 =====================
\section{Bảng tham số cố định}
\begin{center}\small
\begin{tabular}{@{}p{5.2cm} c p{6.5cm}@{}}
\toprule
\textbf{Tham số} & \textbf{Ký hiệu} & \textbf{Giá trị}\\
\midrule
Bước thời gian & $T_s$ & $0.05$ s\\
Chân trời dự báo & $N$ & $20$\\
Chân trời điều khiển & $N_c$ & $5$ (move-blocking)\\
Số kịch bản robust & $M$ & $5$\\
Cửa sổ co rút & $H$ & $20$ bước ($1.0$ s)\\
Trọng số cuối kỳ / delta-u & $Q_f,\,dU$ & $0,\ 0$\\
Độ dài mỗi quỹ đạo & -- & 1000 bước $=50$ s\\
Episode & -- & 5 họ quỹ đạo $\times$ 1000 bước\\
Hạt giống (song song) & seed & $260914001\dots260914005$\\
Surrogate & lr/batch/buffer & $10^{-3}$ / 256 / $10^{5}$\\
SAC & lr/$\gamma$/batch & $3{\times}10^{-4}$ / 0 (bandit) / 256\\
Trọng số khởi tạo & $Q,R$ & Bryson (Mục 5.2)\\
Trọng số residual triển khai & $\alpha$ & $\cS\cdot g_L(\cL)$\\
Ngưỡng cổng LQR & $c_{\mathrm{low}},c_{\mathrm{high}}$ & thiết kế, $0\le c_{\mathrm{low}}<c_{\mathrm{high}}\le1$\\
Thang sai số cho $\cS$ & $\varepsilon_p$ & thiết kế, công bố (dùng trong $s_k$)\\
Cửa sổ RMS của $\cS$ & -- & 20 bước VỪA QUA (past window)\\
\bottomrule
\end{tabular}
\end{center}

% ===================== 14 =====================
\section{Tài liệu tham khảo}
\begin{enumerate}[leftmargin=2.0em,itemsep=1pt,label={[\arabic*]}]
\item Bryson, A. E., \& Ho, Y. C. (1975). \emph{Applied Optimal Control: Optimization, Estimation and
Control}. Hemisphere/Taylor \& Francis. (Quy tắc Bryson.)
\item Bouabdallah, S., Noth, A., \& Siegwart, R. (2004). PID vs LQ control techniques applied to an indoor
micro quadrotor. \emph{IEEE/RSJ IROS 2004}, vol. 3, pp. 2451--2456. DOI 10.1109/IROS.2004.1389776.
\item Sir Elkhatem, A., \& Naci Engin, S. (2022). Robust LQR and LQR-PI control strategies based on
adaptive weighting matrix selection for a UAV position and attitude tracking control. \emph{Alexandria
Engineering Journal}, 61, 6275--6292.
\item Anderson, B. D. O., \& Moore, J. B. (1990). \emph{Optimal Control: Linear Quadratic Methods}.
Prentice-Hall.
\item Khalil, H. K. (2002). \emph{Nonlinear Systems} (3rd ed.). Prentice-Hall.
\item Calafiore, G. C., \& Campi, M. C. (2006). The scenario approach to robust control design. \emph{IEEE
Transactions on Automatic Control}, 51(5), 742--753.
\item Haarnoja, T., Zhou, A., Abbeel, P., \& Levine, S. (2018). Soft Actor-Critic: Off-Policy Maximum
Entropy Deep RL with a Stochastic Actor. \emph{ICML 2018}.
\item Verschueren, R., et al. (2022). acados -- a modular open-source framework for fast embedded optimal
control. \emph{Mathematical Programming Computation}, 14, 147--183.
\item Andersson, J. A. E., Gillis, J., Horn, G., Rawlings, J. B., \& Diehl, M. (2019). CasADi: a software
framework for nonlinear optimization and optimal control. \emph{Mathematical Programming Computation}, 11,
1--36.
\item Ross, S., Gordon, G., \& Bagnell, D. (2011). A Reduction of Imitation Learning and Structured
Prediction to No-Regret Online Learning (DAgger). \emph{AISTATS 2011}.
\item Bishop, C. M. (2006). \emph{Pattern Recognition and Machine Learning}. Springer.
\end{enumerate}

\end{document}
